\documentclass{article}

\usepackage{float}
\usepackage{amssymb}
\usepackage{amsmath}
\usepackage{hyperref}
\usepackage{graphicx} % Required for inserting images
\usepackage{booktabs} % For professional-looking table rules
\usepackage{caption} % The booktabs package will give you nicer line spacing, but if you don't care about the spacing, you can use the standard hline command
\usepackage{subcaption} % For creating subfigures and subtables
\usepackage{float} % The standard floats.
\usepackage{mathrsfs}
\usepackage{footnote}
\usepackage{algorithm}
\usepackage{algorithmic}
\usepackage{bm}
\usepackage{hyperref}
\hypersetup{
    colorlinks=true,
    linkcolor=red,
    filecolor=red,
    urlcolor=red,
    citecolor=red
    }

\urlstyle{same}
\usepackage[left=4cm, right=4cm, top = 3cm, bottom = 3cm]{geometry}


\title{Phase Diversity Adaptive Optics for Multiphoton Microscopy}
\author{ Max Watzky\footnote{Departments of Physics and Mathematics, Yale University}, \ Shengqi Wang\footnote{Department of Biomedical Engineering, Yale University}, \ and Professor Cristina Rodr\'{\i}guez\footnote{Departments of Applied Physics and Biomedical Engineering, Yale University}}
\date{September 21st, 2026}

\begin{document}
\maketitle

\section{Forward Model}
Let $\mathbf{u}$ be the coordinates of the back focal plane (which we treat as two-dimensional), and
$\mathbf{x}$ be the coordinates of the front focal plane (which we treat as three-dimensional). We denote a 3D stack of images as a map from $\mathbb{R}^3$ to the nonnegative real line:
\begin{equation}
d: \mathbb{R}^3 \to [0, \infty)
\end{equation}
The image stack is generated by convolving a 3D PSF $s_k: \mathbb{R}^3 \to [0, \infty)$ with the underlying 3D object $f: \mathbb{R}^3 \to [0, \infty)$. It is often convenient to consider the individual images separately. We will label the images with the index $k$, and write:
\begin{equation}
d_k(x, y) = \sum_{i=1}^P f(x, y, z_i) \otimes s_k(x, y, z_i - z_k)
\end{equation}
where $z_k$ is the focal plane of the image, and $\{z_i\}$ are the $P$ planes around this focal plane.

The PSF is generated by taking the pupil function $H_k: \mathbb{R}^2 \to \mathbb{C}$, applying defocus according to the z-component of $\mathbf{x}$, performing the inverse Fourier transform, and raising the squared magnitude of the result to the $N$th power:
\begin{equation}
s_k(\mathbf{x}) = |h_k(\mathbf{x})|^{2N}  \ \ \ \ h_k(\mathbf{x}) = \mathcal{F}^{-1} \{H_k(\mathbf{u}) e^{i z \gamma(\mathbf{u})}\}
\end{equation}
The extra defocus is what sets this model apart from 2D. We calculate it with:
\begin{equation}
\gamma(\mathbf{u}) = \frac{2\pi n_m}{\lambda} \sqrt{1-  \left(\frac{\lambda \|\mathbf{u}\|}{n_m} \right)^2}
\end{equation}
where $n_m$ is the refractive index and $\lambda$ is the vacuum wavelength.
Furthermore, the pupil function is generated by a pupil mask/apodization $P(\mathbf{u})$, an underlying phase aberration $\phi(\mathbf{u})$ and a known diversity phase $\theta_k(\mathbf{u})$. In all, we write this as:
\begin{equation}
H_k(\mathbf{u}) = |P(\mathbf{u})| e^{i \left(\phi(\mathbf{u}) +  \theta_k(\mathbf{u})\right)}
\end{equation}
Sometimes, we will expand $\phi(\mathbf{u})$ in a basis $\{ Z_i\}$. In this case, we will use the vector $\bm{\alpha}$ to describe the coefficients of $\phi$ in this basis.

\section{Cost Functionals}
Let $\{d_k\}$ be the set of acquired images, with known diversity phases $\{\theta_k\}$. The goal is to estimate the underlying aberration $\phi(\mathbf{u})$ and the underlying sample $f(\mathbf{x})$ to minimize a cost functional $\mathcal{L} \{f(\mathbf{x}), \bm{\alpha}\}$. In keeping with the literature, we will treat two cases: Gaussian and Poisson noise. In the Gaussian case, we seek to minimize the least-squares objective:
\begin{equation}
\mathcal{L}_g \{f, \bm{\alpha}\} = \sum_{k=1}^K \sum_{(x,y)} \left[d_k(x, y) - \mu_k(x,y) \right]^2
\label{eq:gaussian_cost_functional}
\end{equation}
where we use the shorthand:
\begin{align}
\mu_k(x, y) = \sum_{i=1}^P f(x, y, z_i) \otimes s_k(x, y, z_i - z_k)
\end{align}
where $s_k$ here denotes the prediction of the PSF from $\bm{\alpha}$. We can also write out the Poisson objective:
\begin{equation}
\mathcal{L}_p \{f, \bm{\alpha}\} = 2 \sum_{k=1}^K \sum_{(x,y)} \left[d_k(x, y) \ln(d_k(x,y)) - d_k(x,y) \ln(\mu_k(x,y)) - (d_k(x,y) - \mu_k(x,y))\right]
\end{equation}

\section{Gaussian Noise}
We will first attempt to solve the problem for the case of Gaussian noise.

\subsection{Object Estimation}
We begin by estimating the underlying object $f$. By Parseval's theorem, we can rewrite the Gaussian cost functional in \autoref{eq:gaussian_cost_functional} in the Fourier domain:
\begin{equation}
\mathcal{L}_g \{F, \bm{\alpha}\} = \sum_{k=1}^K \sum_{(u_x, u_y)} \left| D_k - \sum_{i=1}^P F_{z_i} S_k (z_i - z_k)\right|^2
\label{eq:gaussian_cost_functional_fourier}
\end{equation}
where $D_k$ denotes the Fourier transform of $d_k$, $F_{z_i}$ denotes the Fourier transform of $f(x, y, z_i)$ and $S_k(z_i - z_k)$ denotes the Fourier transform of the $k$th simulated PSF slice at $z_i - z_k$. The notation $\mathcal{F}_{xy}$ used below denotes the same lateral Fourier transform.
At each fixed spatial frequency, the unconstrained object that minimizes the cost functional will satisfy $\partial \mathcal{L}_g/ \partial F_{z_j}^* = 0$ for each $j=1,\ldots,P$. We compute each of these w.r.t the complex conjugate of $F_{z_j}$ so that the resulting expression will include $F_{z_j}$ directly:
\begin{align}
\frac{\partial \mathcal{L}_g}{\partial F_{z_j}^*} &= \frac{\partial}{\partial F_{z_j}^*} \sum_{k=1}^K \sum_{(u_x, u_y)} \left| D_k - \sum_{i=1}^P F_{z_i} S_k (z_i - z_k)\right|^2 \\
&=  \sum_{k=1}^K \frac{\partial}{\partial F_{z_j}^*} \left[ \left(D_k - \sum_{i=1}^P F_{z_i} S_k (z_i - z_k)\right) \left(D^*_k - \sum_{i=1}^P F^*_{z_i} S^*_k (z_i - z_k)\right) \right] \\
&=  - \sum_{k=1}^K \left[ S_k^*(z_j- z_k) \left(D_k - \sum_{i=1}^P F_{z_i} S_k (z_i - z_k)\right) \right]
\end{align}
Rearranging, we get:
\begin{align}
\sum_{i=1}^P \left[\sum_{k=1}^K S_k^*(z_j - z_k) S_k(z_i - z_k)  \right] F_{z_i} = \sum_{k=1}^K S_k^*(z_j - z_k) D_k
\label{eq:optimize_for_f}
\end{align}
We get $P$ equations of this form, one for each index $j=1,...,P$. Together, they form the matrix equation:
\begin{align}
(\mathbf{S}^\dagger \mathbf{S}) \mathbf{F} = \mathbf{S}^\dagger \mathbf{D}
\end{align}
It is helpful to briefly remind ourselves of the dimensions of these matrices: $\mathbf{F}$ is a single column with one row for each equation ($P$). By similar logic, we can infer that $\mathbf{D}$ has one column and one row for each image ($K$). $\mathbf{S}$ is $K \times P$, and $\mathbf{S}^\dagger$ is $P \times K$, so, $\mathbf{S}^\dagger \mathbf{S}$ is $P \times P$.
We can thus verify the dimensions of our equation are correct. Solving, we get the estimator:
\begin{align}
\hat{\mathbf{F}} = (\mathbf{S}^\dagger \mathbf{S})^{-1} \mathbf{S}^\dagger \mathbf{D}
\end{align}
where we have assumed that $\mathbf{S}^\dagger \mathbf{S}$ is invertible. This means that all of the columns of $\mathbf{S}$ must be linearly independent. Since $\mathbf{S}$ is $K \times P$, it must be the case that $K \ge P$; otherwise this is impossible. As Kner points out, this tells us that we must have at least $P$ measurements to estimate the image at $P$ planes.

Sometimes, we would like an estimate, even when we have insufficient data. We therefore introduce a fixed regularization parameter $\rho>0$ as follows:
\begin{align}
\hat{\mathbf{F}} = (\mathbf{S}^\dagger \mathbf{S} + \rho I_{P})^{-1} \mathbf{S}^\dagger \mathbf{D}
\label{eq:regularized_F}
\end{align}
which will allow us to solve for the inverse even if $\mathbf{S}$ has insufficient rank. This is equivalent to placing a squared $L^2$ norm penalty on $\mathbf{F}$ in the original cost function \autoref{eq:gaussian_cost_functional_fourier}.

\subsection{Aberration Estimation}
Next, we will discuss the process of estimating the aberration along a basis $\{Z_i\}$, with real coefficients $\alpha_i$:
\begin{equation}
\phi(\mathbf{u}) = \sum_i \alpha_i Z_i(\mathbf{u})
\end{equation}
We will begin by defining the following variable, the residual between our $k$th estimated image (which we generate from the estimated object $\hat{\mathbf{F}}$) and the actual image $D_k$:
\begin{equation}
r_k = D_k - \sum_{i=1}^P \hat{F}_{z_i} S_{k}(z_i - z_k)
\end{equation}
With $\hat{\mathbf{F}}$ plugged in, and the additional squared $L^2$ penalty added, we rewrite the original cost functional \autoref{eq:gaussian_cost_functional_fourier} as:
\begin{equation}
\mathcal{L}_g \{\bm{\alpha}\} = \sum_{(u_x, u_y)} \sum_{k=1}^K r_k r_k^*  + \rho \sum_{(u_x, u_y)} \sum_{i=1}^P \hat{F}_{z_i} \hat{F}^*_{z_i}
\end{equation}
To take the gradient $\nabla \mathcal{L}_g$, we begin by taking the derivative with respect to an arbitrary aberration coefficient $\alpha_m$:
\begin{align}
\frac{\partial \mathcal{L}_g}{\partial \alpha_m} &= \sum_{(u_x, u_y)} \sum_{k=1}^K \left[r_k^* \left(\frac{\partial r_k}{\partial \alpha_m} \right)  + r_k \left(\frac{\partial r_k}{\partial \alpha_m} \right)^*\right] + \rho \sum_{(u_x, u_y)} \sum_{i=1}^P \left[\hat{F}_{z_i}^* \left(\frac{\partial \hat{F}_{z_i}}{\partial \alpha_m} \right)  + \hat{F}_{z_i} \left(\frac{\partial \hat{F}_{z_i}}{\partial \alpha_m} \right)^*\right]
\end{align}
In both terms, we recognize sums of the form $w + w^* = 2 \Re(w)$. Using this identity, we rewrite:
\begin{equation}
\frac{\partial \mathcal{L}_g}{\partial \alpha_m} = 2 \Re \left[\sum_{(u_x, u_y)} \sum_{k=1}^K r_k^* \frac{\partial r_k}{\partial \alpha_m} \right] + 2 \rho \Re \left[\sum_{(u_x, u_y)} \sum_{i=1}^P \hat{F}_{z_i}^* \frac{\partial \hat{F}_{z_i}}{\partial \alpha_m} \right]
\end{equation}
We examine the next derivative:
\begin{equation}
\frac{\partial r_k}{\partial \alpha_m} = \frac{\partial}{\partial \alpha_m} \left[D_k - \sum_{i=1}^P \hat{F}_{z_i} S_k(z_i - z_k) \right] = - \sum_{i=1}^P \left[\frac{\partial \hat{F}_{z_i}}{\partial \alpha_m} S_k(z_i - z_k) + \hat{F}_{z_i} \frac{\partial S_k(z_i - z_k)}{\partial \alpha_m} \right]
\end{equation}
From here on, for compactness we denote:
\begin{align}
S_{ki} = S_k(z_i - z_k)
\end{align}
Plugging in our results, we get:
\begin{align}
\frac{\partial \mathcal{L}_g}{\partial \alpha_m} = -2 \Re \left[\sum_{(u_x, u_y), k, i} r_k^* S_{ki} \frac{\partial \hat{F}_{z_i}}{\partial \alpha_m}\right] - 2 \Re \left[ \sum_{(u_x, u_y), k, i} r_k^* \hat{F}_{z_i} \frac{\partial S_{ki}}{\partial \alpha_m} \right] + 2 \rho \Re \left[\sum_{(u_x, u_y), i} \hat{F}^*_{z_i} \frac{\partial \hat{F}_{z_i}}{\partial \alpha_m}\right]
\end{align}
We briefly examine the first and third terms, since they both involve the derivative of the optimized object with respect to the aberration coefficient. Inside the real-part operator, we have:
\begin{align}
-\sum_{(u_x, u_y)} \sum_i \left[\sum_k r_k^* S_{ki} \right] \frac{\partial \hat{F}_{z_i}}{\partial \alpha_m} + \rho \sum_{(u_x, u_y)} \sum_i \hat{F}^*_{z_i} \frac{\partial \hat{F}_{z_i}}{\partial \alpha_m}
\end{align}
However, recall from \autoref{eq:regularized_F} that:
\begin{align}
(\mathbf{S}^\dagger \mathbf{S} + \rho I_P) \hat{\mathbf{F}} = \mathbf{S}^\dagger (\mathbf{r} + \mathbf{S} \hat{\mathbf{F}})
\end{align}
which gives:
\begin{align}
\rho \hat{\mathbf{F}} = \mathbf{S}^\dagger \mathbf{r}
\end{align}
which allows us to cancel the corresponding terms in the summation. The result is:
\begin{align}
\frac{\partial \mathcal{L}_g}{\partial \alpha_m} = - 2 \Re \left[ \sum_{(u_x, u_y), k, i} r_k^* \hat{F}_{z_i} \frac{\partial S_{ki}}{\partial \alpha_m} \right]
\end{align}
The next task is to evaluate the PSF derivative. Recall the forward model:
\begin{align}
\frac{\partial S_{ki}}{\partial \alpha_m} = \frac{\partial}{\partial \alpha_m} \mathcal{F} \{s_k(x, y, z_i - z_k) \} = \mathcal{F} \left\{ \frac{\partial s_k(x, y, z_i - z_k)}{\partial \alpha_m} \right\}
\end{align}
where
\begin{align}
\frac{\partial s_k(x, y, z_i - z_k)}{\partial \alpha_m} = \frac{\partial}{\partial \alpha_m} (h_k h_k^*)^N = 2 N|h_k|^{2N-2} \Re \left[h_k^* \frac{\partial h_k}{\partial \alpha_m} \right]
\end{align}
where
\begin{align}
\frac{\partial h_k}{\partial \alpha_m} = \frac{\partial}{\partial \alpha_m} \mathcal{F}^{-1} \{H_k(\mathbf{u}) e^{i z \gamma(\mathbf{u})}\} = \mathcal{F}^{-1} \left\{i Z_m(\mathbf{u}) H_k(\mathbf{u}) e^{iz \gamma(\mathbf{u})} \right\}
\end{align}
where we have used that the only $\alpha_m$ dependence in this term is in $H_k$, where:
\begin{align}
\frac{\partial H_k}{\partial \alpha_m} = \frac{\partial }{\partial \alpha_m} |P(\mathbf{u})| e^{i \left(\phi(\mathbf{u}) +  \theta_k(\mathbf{u})\right)} = i Z_m H_k
\end{align}
by expanding $\phi$ in our basis as $\sum_j Z_j \alpha_j$. Summarizing, we have:
\begin{equation}
\begin{aligned}
\frac{\partial \mathcal{L}_g}{\partial \alpha_m} ={}& -4N\,\Re \sum_{(u_x,u_y)} \sum_{k=1}^{K} \left[D_k-\sum_{i=1}^{P}\hat{F}_{z_i}S_k(z_i-z_k)\right]^*
\\&\quad\times\left[\sum_{j=1}^{P}\hat{F}_{z_j}\,\mathcal{F}\left\{\left|h_k(z_j-z_k)\right|^{2N-2}\Re\left[h_k^*(z_j-z_k)\,\mathcal{F}^{-1}\left\{iZ_mH_k e^{i(z_j-z_k)\gamma}\right\}\right]\right\}\right].
\end{aligned}
\end{equation}
I then used GPT-6 Astra to automatically compute the Hessian, which takes the following form:
\begin{equation}
\begin{aligned}
\frac{\partial^2 \mathcal{L}_g}
{\partial\alpha_m\,\partial\alpha_\ell}
={}&
8N^2\operatorname{Re}
\sum_{(u_x,u_y)}\sum_{k=1}^{K}
\left[
\sum_{i=1}^{P}\widehat F_{z_i}\,
\mathcal F_{xy}
\left\{
|h_k(z_i-z_k)|^{2N-2}a_{ki,m}
\right\}
\right]^*
\\
&\quad\times
\left[
\sum_{j=1}^{P}\widehat F_{z_j}\,
\mathcal F_{xy}
\left\{
|h_k(z_j-z_k)|^{2N-2}a_{kj,\ell}
\right\}
\right]
\\[4pt]
&-
8N(N-1)\operatorname{Re}
\sum_{(u_x,u_y)}\sum_{k=1}^{K}
r_k^*
\sum_{j=1}^{P}\widehat F_{z_j}\,
\mathcal F_{xy}
\left\{
|h_k(z_j-z_k)|^{2N-4}a_{kj,m}a_{kj,\ell}
\right\}
\\[4pt]
&-
4N\operatorname{Re}
\sum_{(u_x,u_y)}\sum_{k=1}^{K}
r_k^*
\sum_{j=1}^{P}\widehat F_{z_j}\,
\mathcal F_{xy}
\left\{
|h_k(z_j-z_k)|^{2N-2}b_{kj,m\ell}
\right\}
\\[4pt]
&-
2\operatorname{Re}
\sum_{(u_x,u_y)}
\sum_{i=1}^{P}\sum_{j=1}^{P}
c_{i,m}^*
\left[
\left(\mathbf S^\dagger\mathbf S+\rho I_P\right)^{-1}
\right]_{ij}
c_{j,\ell}.
\end{aligned}
\end{equation}
where:
\begin{align*}
a_{ki,m}
&=
\operatorname{Re}\left[
h_k^*(z_i-z_k)\,
\mathcal F^{-1}
\left\{
iZ_mH_k e^{i(z_i-z_k)\gamma}
\right\}
\right],
\\[6pt]
b_{ki,m\ell}
&=
\operatorname{Re}\Bigg[
\left(
\mathcal F^{-1}
\left\{
Z_\ell H_k e^{i(z_i-z_k)\gamma}
\right\}
\right)^*
\mathcal F^{-1}
\left\{
Z_mH_k e^{i(z_i-z_k)\gamma}
\right\}
\\
&\qquad\qquad-
h_k^*(z_i-z_k)\,
\mathcal F^{-1}
\left\{
Z_mZ_\ell H_k e^{i(z_i-z_k)\gamma}
\right\}
\Bigg],
\\[6pt]
c_{j,m}
&=
2N\sum_{k=1}^{K}\Bigg[
\left(
\mathcal F_{xy}
\left\{
|h_k(z_j-z_k)|^{2N-2}a_{kj,m}
\right\}
\right)^*r_k
\\
&\qquad\qquad-
S_k^*(z_j-z_k)
\sum_{i=1}^{P}\widehat F_{z_i}\,
\mathcal F_{xy}
\left\{
|h_k(z_i-z_k)|^{2N-2}a_{ki,m}
\right\}
\Bigg],
\\[6pt]
r_k
&=
D_k-\sum_{i=1}^{P}\widehat F_{z_i}S_k(z_i-z_k),
\\[6pt]
h_k(z)
&=
\mathcal F^{-1}\left\{H_k e^{iz\gamma}\right\},
\\[6pt]
S_k(z)
&=
\mathcal F_{xy}\left\{|h_k(z)|^{2N}\right\},
\\[6pt]
[\mathbf S]_{ki}
&=
S_k(z_i-z_k),
\\[6pt]
\mathbf D
&=
(D_1,\ldots,D_K)^{\mathsf T},
\\[6pt]
\widehat{\mathbf F}
&=
(\widehat F_{z_1},\ldots,\widehat F_{z_P})^{\mathsf T}
=
\left(\mathbf S^\dagger\mathbf S+\rho I_P\right)^{-1}
\mathbf S^\dagger\mathbf D,
\\[6pt]
\mathcal{L}_g(\bm{\alpha})
&=
\sum_{(u_x,u_y)}
\left[
\sum_{k=1}^{K}|r_k|^2
+\rho\sum_{i=1}^{P}|\widehat F_{z_i}|^2
\right].
\end{align*}

\begin{align}
\Delta \mathbf{\alpha} = -\mathbf{H}^{-1} \mathbf{g}
\end{align}
\end{document}